\documentclass[11pt,a4paper]{ctexart}
\usepackage[a4paper,margin=1.60cm,top=1.8cm,bottom=1.8cm]{geometry}
\usepackage{../../../00_system/templates/ipara-reading}

% ==============================================================================
% 《English Decoded!》单篇精读训练手札
% 原文与首读痕迹：article.md
% 本文件：完整呈现本篇原文，并保存首读断点、结构、声音参考、真实输出与对照。
% 不要求学习者在同一时段继续机械复述；后续只做间隔后的冷检索。
% ==============================================================================

\begin{document}

\articleheader
  {English Decoded!}
  {What is said -- and what is meant}
  {Spotlight · Language}
  {Pragmatics · Communication Styles · Hidden Meaning}
  {精读 · 语用判断 · 口语重构}
  {2026-08-20}

% ------------------------------------------------------------------------------
% 一、双栏精读：完整原文 + 当前真正需要的点拨
% ------------------------------------------------------------------------------
\readingsection{一、双栏精读}

\begin{paracol}{2}

\artpara{P1}{
  LEARNING A FOREIGN LANGUAGE is a lot of fun but also quite challenging -- \obs{not helped by the fact that}{1} people communicate in very different ways. Millions of people across the world speak English, so providing rules about how people communicate specifically in English is difficult. \obs{That said}{2}, English, like every language, has nuanced ways of sending certain messages. And that is our focus here.
}
\switchcolumn
\noteobs{1}{not helped by the fact that...}
  {不是“没有得到帮助”，而是“这个事实让原本的困难变得更严重”。可口语重构为 \textit{...and the fact that people communicate differently makes it even harder}.}
\noteobs{2}{That said, ...}
  {先承认前面的限制，再转入仍然成立的观点，接近“话虽如此 / 尽管如此”。这里的逻辑是：英语使用者差异很大，很难下绝对规则；尽管如此，英语仍有可学习的语用模式。}

\switchcolumn*
\artpara{P2}{
  We'll help you appreciate the difference between what people say and what they actually mean in English -- in many different situations. First, look at the box on page 13 to learn about three types of communicators. Then, move on to the quiz to see examples of what each type might typically say in English in a certain scenario -- and, more importantly, to find out how well you can \kw{read between the lines}{3}.
}
\switchcolumn
\noteitem{3}{read between the lines · 本篇核心表达}
  {不只按字面理解，而是从语气、语境和措辞中读出隐藏意图。}
  {understand what someone really means, not just the literal words}
  {学习者第二次复述说成 \textit{read the air}；本文应保留原表达 \textit{read between the lines}，后续只在间隔检索时再检查，不要求今天继续重复。}

\switchcolumn*
\artpara{P3}{
  \kw{Assertive communicators}{4} have mastered the art of being ``firm but fair''. They stand up for their rights and set clear boundaries, while still being respectful of other people. Assertive communication tends to be pleasant and straightforward.
}
\switchcolumn
\noteitem{4}{assertive · firm but fair}
  {明确表达自己的需求、权利和边界，同时尊重别人。它不是 aggressive 的温和版，而是一种“清晰 + 尊重”的沟通方式。}
  {direct and respectful}
  {学习者首读 [?]；第二次输出已经能补回“直接表达 + 尊重别人”的边界。}

\switchcolumn*
\artpara{P4}{
  Aggressive communicators are straight talkers, who let people know exactly how they feel. They can sound aggressive and often cause unnecessary upset -- but at least they \kw{get their point across}{5} and \kw{you'll know where you stand with them}{6}.
}
\switchcolumn
\noteitem{5}{get one's point across}
  {把自己的意思明确传达到，让对方听懂。这里强调 aggressive 沟通虽然伤人，但意图通常不隐藏。}
  {make your meaning clear}
  {当前留在本篇。}
\noteitem{6}{know where you stand with someone}
  {知道对方真实态度和立场，不需要猜。这里和 passive-aggressive 的“hidden message”形成对照。}
  {know exactly how they feel about you / the situation}
  {当前留在本篇。}

\switchcolumn*
\artpara{P5}{
  Passive-aggressive communicators avoid direct conflict. They use sarcasm to express their feelings or say the opposite of what they mean -- a style of communication that only works, of course, if the other person gets the hidden message.
}
\switchcolumn
\noteitem{7}{passive-aggressive · 核心边界}
  {重点不是“= sarcasm”，而是\textbf{避免直接冲突 + 用间接方式传递隐藏负面信息}。sarcasm 和反话只是常见手段。}
  {indirect criticism / hidden annoyance}
  {这是本篇最需要和 assertive、aggressive 区分开的概念。}

\switchcolumn*
\artpara{P6}{
  Would you be able to recognize, let's say, the sarcasm in someone's seemingly helpful response? Would you fall for a friendly offer that's actually not an offer at all but a reproach? Would you realize that your conversation partner is, in fact, not as understanding as he or she sounds?
}
\switchcolumn
\noteobs{8}{seemingly helpful / actually a reproach}
  {这一段直接点明 quiz 的难点：真正危险的不是明显骂人的句子，而是\textbf{表面礼貌、实际在责怪或讽刺}的句子。}

\switchcolumn*
\artpara{P7}{
  Here are ten everyday situations and a choice of three responses for each one. Which response should \kw{put you on your guard}{9} -- and best not be \kw{taken literally}{10}? Take a moment to consider what your conversation partner is really saying, then tick the box.
}
\switchcolumn
\noteitem{9}{put someone on their guard}
  {让某人提高警惕、意识到这里可能有隐藏风险。}
  {make you cautious / alert}
  {本篇语境中就是：哪句话不能只按表面礼貌程度判断。}
\noteitem{10}{take something literally}
  {按字面意思理解。本文训练恰恰要求识别哪些句子不能这样读。}
  {understand the exact surface meaning only}
  {与 \textit{read between the lines} 构成一组对照。}

\switchcolumn*
\artpara{P8--P9}{
  \textbf{1. A colleague is already struggling with their own work.} You ask a colleague for help, but they are already struggling with their own work and say:
  \begin{itemize}[leftmargin=*]
    \item A. ``That's not my responsibility and I'm busy.''
    \item B. ``I can't do anything today, but if you make a start, I could find some time to help you on Friday, if you still need it.''
    \item C. ``Sure. I guess I could do your work as well as my own. Why not?''
  \end{itemize}
}
\switchcolumn
\noteobs{11}{Quiz 1 · hidden response = C}
  {A 是直接拒绝；B 是 assertive / cooperative；C 表面同意，实际是在讽刺“你是不是想让我把你的活也全做了”。}

\switchcolumn*
\artpara{P10--P11}{
  \textbf{2. Dirty dishes in the sink.} You have left a pile of dirty dishes in the sink. Your housemate says:
  \begin{itemize}[leftmargin=*]
    \item A. ``I'll just put my rubber gloves on and wash these few dishes. It's not a big deal, is it?''
    \item B. ``I've noticed there are some dirty dishes in the sink. Could you wash them by bedtime, so everything is clean for the end of the day?''
    \item C. ``You need to wash these dishes.''
  \end{itemize}
}
\switchcolumn
\noteobs{12}{Quiz 2 · hidden response = A}
  {学习者选对了 A，但最初写成 \textit{direct to criticize}。它恰恰不是 direct：说话者自己去洗，并用 ``It's not a big deal, is it?'' 间接表达不满。B 是较典型的 assertive request。}

\switchcolumn*
\artpara{P12--P13}{
  \textbf{3. An upcoming deadline.} You remind a staff member of an upcoming deadline. They say:
  \begin{itemize}[leftmargin=*]
    \item A. ``Don't worry. I have it on my schedule for Thursday. I'll let you know if I hit any problems, otherwise, \kw{it's all in hand}{13}.''
    \item B. ``I know. It's in my calendar. I don't need you to micromanage me.''
    \item C. ``Oh, I'm sorry I haven't already sent the work in ahead of the deadline. Would you like me to send you daily updates?''
  \end{itemize}
}
\switchcolumn
\noteitem{13}{it's all in hand}
  {事情已经安排好、在掌控中，接近 \textit{under control}。学习者首读 [?]。}
  {Don't worry -- I've got it under control.}
  {A 是正常 reassurance；真正的 hidden response 是 C。}
\noteobs{14}{Quiz 3 · hidden response = C}
  {C 表面道歉并主动提出“每天汇报”，实际在讽刺对方过度催促。这个题说明：越礼貌不一定越友好。}

\switchcolumn*
\artpara{P14--P15}{
  \textbf{4. Jumping the queue.} You're in a terrible rush and give in to the temptation to jump the queue. Someone says:
  \begin{itemize}[leftmargin=*]
    \item A. ``Excuse me! There's a queue here.''
    \item B. ``No, please, go right ahead! After you!''
    \item C. ``Hey! You're pushing in! Get to the back of the queue like everyone else!''
  \end{itemize}
}
\switchcolumn
\noteobs{15}{Quiz 4 · hidden response = B}
  {B 的表面意思是“请先走”，实际意思相反：你插队让我很恼火。学习者首读判断正确。}

\switchcolumn*
\artpara{P16--P17}{
  \textbf{5. Unasked-for career advice.} With the best intentions, you \kw{volunteer career advice}{16} to a much younger family member. They say:
  \begin{itemize}[leftmargin=*]
    \item A. ``Thank you for your opinion. I'll be sure to bear it in mind.''
    \item B. ``Your advice isn't helpful because you're out of touch.''
    \item C. ``I appreciate that you're looking out for me, but things have changed in the last few years and I'm confident in my decisions.''
  \end{itemize}
}
\switchcolumn
\noteitem{16}{volunteer advice}
  {这里的 \textit{volunteer} 是“主动提供”，尤其指对方没有要求时仍给出建议。学习者首读 [!]。}
  {offer advice without being asked}
  {这个细节决定了后面 A 的“Thank you for your opinion”可能带敷衍意味。}
\noteobs{17}{Quiz 5 · hidden response = A}
  {C 才是很好的 assertive refusal：先承认善意，再明确边界。A 表面礼貌，但在该语境中更像“我听到了，但我并不需要你的建议”。}

\switchcolumn*
\artpara{P18--P19}{
  \textbf{6. A dog barking late at night.} Your dog is barking late at night, keeping everyone awake. A neighbour says:
  \begin{itemize}[leftmargin=*]
    \item A. ``Shut that dog up or I'm calling the police!''
    \item B. ``I'm worried for the wellbeing of the dog -- and its owner. Is something wrong?''
    \item C. ``I wanted to talk to you about your dog. Is there something you can do to settle him down at night?''
  \end{itemize}
}
\switchcolumn
\noteobs{18}{Quiz 6 · hidden response = B}
  {B 用“担心狗和主人”包装了对噪音的抱怨；C 才是直接但尊重的请求。学习者首读判断正确。}

\switchcolumn*
\artpara{P20--P21}{
  \textbf{7. Interrupting a colleague in a meeting.} In a meeting, you have a brilliant idea and can't stop yourself from interrupting the colleague who's currently speaking. She says:
  \begin{itemize}[leftmargin=*]
    \item A. ``I'll just be quiet -- you clearly have something important to say.''
    \item B. ``Will you shut up and let me finish for once?''
    \item C. ``Hang on a second, I'd like to finish my thought before we move on to your point.''
  \end{itemize}
}
\switchcolumn
\noteobs{19}{Quiz 7 · hidden response = A}
  {A 的字面是“那我不说了，你显然更重要”，真实意图是批评你打断别人。C 是典型 assertive boundary。}

\switchcolumn*
\artpara{P22--P23}{
  \textbf{8. Cancelling plans repeatedly.} You have had to cancel plans with a friend a few times, sometimes at the very last minute. He says:
  \begin{itemize}[leftmargin=*]
    \item A. ``You're very unreliable. I don't want to make any more plans with you.''
    \item B. ``I'm sure we'll eventually manage to find a date you can do, if your busy schedule allows.''
    \item C. ``I've noticed we've had to cancel a few things lately. I'd really like to see you, but it feels rather one-sided at the moment.''
  \end{itemize}
}
\switchcolumn
\noteobs{20}{Quiz 8 · hidden response = B}
  {B 表面体谅“你太忙”，实际是在讽刺你总是取消。C 反而把受伤感受和边界说得更直接。}

\switchcolumn*
\artpara{P24--P25}{
  \textbf{9. Looking at your phone while your partner talks.} You are reading the news on your phone while your partner is trying to tell you about their day. Your partner says:
  \begin{itemize}[leftmargin=*]
    \item A. ``Put that phone down! You never listen to a word I say.''
    \item B. ``I'd really like to chat. Can you put the phone down?''
    \item C. ``I'll call my mum instead. She's always interested in what I've been doing.''
  \end{itemize}
}
\switchcolumn
\noteobs{21}{Quiz 9 · hidden response = C}
  {不是 straight，而是 indirect reproach：不直接说“你没在听”，改用“我去找我妈聊，她会听我说”来传递不满。}

\switchcolumn*
\artpara{P26--P27}{
  \textbf{10. Not thanking someone who holds the door.} Someone holds the door open for you and, completely lost in thought, you don't thank them. They say:
  \begin{itemize}[leftmargin=*]
    \item A. ``You're welcome.''
    \item B. ``That was rude!''
    \item C. Nothing. Life's too short...
  \end{itemize}
}
\switchcolumn
\noteobs{22}{Quiz 10 · hidden response = A}
  {A 在没有听到 ``thank you'' 的情况下主动说 ``You're welcome''，形成明显讽刺。C 不是 sarcasm，而是“不值得为这点小事计较”。}

\switchcolumn*
\artpara{P28--P29}{
  Here are some classic passive-aggressive phrases -- complete with a guide to their (not so) secret hidden meanings:
  \begin{itemize}[leftmargin=*]
    \item ``As you are aware, ...'' = I suspect you haven't got a clue what I'm talking about.
    \item ``For future reference, ...'' = Don't make the same silly mistake again!
    \item ``As I stated in my last email, ...'' = You either didn't read or didn't understand my last email.
    \item ``A gentle reminder that...'' = You never stick to your deadlines!
    \item ``Don't worry about me!'' = You should soooo worry -- I need your undivided attention right now!
    \item ``I guess I'll just do it myself.'' = You're not going to do it anyway -- or not well enough. In fact, you never help me!
    \item ``I have nothing better to do with my time.'' = I'm incredibly busy, as usual, but no one cares, as usual.
    \item ``Better late than never!'' = How difficult can it be to be on time? You're always late!
    \item ``I felt a bit silly.'' = I was made to feel a complete idiot because of you! You'd better apologize -- and \kw{a bunch of flowers won't be enough}{23}!
    \item ``It's not your fault, really!'' = Of course it's your fault! Who else's fault would it be?
    \item ``My mistake entirely!'' = You got it wrong! How could you possibly make such a dumb mistake?
    \item ``This really is quite tricky!'' = How can you not understand something so simple?
  \end{itemize}
}
\switchcolumn
\noteitem{23}{a bunch of flowers won't be enough}
  {\textit{a bunch of flowers} 是“一束花”。这里故意把表面轻描淡写的 \textit{I felt a bit silly} 解码成非常强烈的受辱感和责怪。学习者首读 [?]。}
  {You really embarrassed me, and a simple apology won't be enough.}
  {这一整组例子都在展示“surface wording 轻，hidden emotion 重”的落差。}
\noteobs{24}{整页的真正规律}
  {被动攻击并不依赖某一个固定词，而依赖\textbf{字面意义、语境和真实意图之间的不一致}。因此不要背“哪句话一定 passive-aggressive”，而要学会结合情境判断。}

\end{paracol}

% ------------------------------------------------------------------------------
% 二、文章结构与主旨复原
% ------------------------------------------------------------------------------
\readingsection{二、文章结构与主旨复原}

\begin{coretakeaway}[训练后主旨]
  English communication can carry meanings that are not fully expressed by the literal words. The article teaches readers to distinguish among assertive, aggressive, and passive-aggressive communication and to read between the lines in everyday situations.
\end{coretakeaway}

\begin{itemize}
  \item \textbf{问题提出}：英语由全球不同人群使用，因此很难给出绝对沟通规则；但英语同样存在细腻的语用模式。
  \item \textbf{核心分类}：assertive = direct + respectful；aggressive = direct + potentially hurtful；passive-aggressive = indirect + hidden negative message。
  \item \textbf{情境检验}：十个日常场景让读者区分字面内容与真实意图，尤其识别“看起来礼貌、实际上在讽刺或责怪”的句子。
  \item \textbf{总结强化}：最后集中列出常见 passive-aggressive phrases，把“what is said”与“what is meant”并排展示。
\end{itemize}

\begin{coretakeaway}[本篇最重要的概念边界]
  \textbf{Assertive is not weak aggression, and passive-aggressive is not simply sarcasm.} Assertive communication states needs and boundaries clearly while respecting others; aggressive communication is direct but may be hostile; passive-aggressive communication hides criticism or annoyance behind indirect wording.
\end{coretakeaway}

% ------------------------------------------------------------------------------
% 三、语音与韵律训练：只给少量高价值句，不要求今天继续反复练
% ------------------------------------------------------------------------------
\readingsection{三、语音与韵律训练}

\begin{prosodybox}[黄金句 1 · 核心定义]
  \prosodysource{本篇当前未保存可直接播放的真人原音；以下仅作文本级中性朗读参考，不据此评价学习者真实发音。}
  \prosodyorig{Assertive communicators have mastered the art of being ``firm but fair''.}
  \prosodyipa{\textbf{assertive} \ipa{/əˈsɜːtɪv/} \quad$\bm{\cdot}$\quad \textbf{communicator} \ipa{/kəˈmjuːnɪkeɪtə/}}
  \prosodychunks{\stress{ASSERTIVE} communicators \chunk have mastered the art of being \stress{FIRM BUT FAIR}. \tonefall}
  \prosodyfocus{assertive, firm but fair}
  \prosodyflaws{
    \item 当前没有用户原始录音，因此这里只提示词重音和意群，不写“已读对 / 发音错误”。
    \item \textit{firm but fair} 作为一个整体概念说出，不把三个词切散。
  }
\end{prosodybox}

\begin{prosodybox}[黄金句 2 · 文章任务]
  \prosodysource{仅文本参考。}
  \prosodyorig{Take a moment to consider what your conversation partner is really saying.}
  \prosodyipa{\textbf{consider} \ipa{/kənˈsɪdə/} \quad$\bm{\cdot}$\quad \textbf{conversation} \ipa{/ˌkɒnvəˈseɪʃən/}}
  \prosodychunks{Take a \stress{MOMENT} \chunk to consider \chunk what your conversation partner is \stress{REALLY SAYING}. \tonefall}
  \prosodyfocus{moment, really saying}
  \prosodyflaws{
    \item 训练重点是把 \textit{what your conversation partner is really saying} 当作完整信息块，而不是逐词拼接。
  }
\end{prosodybox}

\begin{prosodybox}[黄金句 3 · 可迁移表达]
  \prosodysource{仅文本参考。}
  \prosodyorig{The article teaches readers to distinguish among assertive, aggressive, and passive-aggressive communication.}
  \prosodyipa{\textbf{distinguish} \ipa{/dɪˈstɪŋɡwɪʃ/} \quad$\bm{\cdot}$\quad \textbf{aggressive} \ipa{/əˈɡresɪv/}}
  \prosodychunks{The article teaches readers \chunk to \stress{DISTINGUISH AMONG} \chunk assertive, aggressive, and passive-aggressive communication. \tonefall}
  \prosodyfocus{distinguish among; three-way contrast}
  \prosodyflaws{
    \item 这句只作为以后间隔检索时的声音模板；今天不要求继续重复整篇。
  }
\end{prosodybox}

% ------------------------------------------------------------------------------
% 四、真实输出、参考版与本轮收束
% ------------------------------------------------------------------------------
\readingsection{四、真实输出与对照}

\oralblock{第一次冷输出 · 原始转写}{
  There is many nuance In English while we chat with others and this article maybe trying to teach me how to distinct these Ative reply is firm and direct, aggressive, it's unkind and Ative reply is firm and direct, aggressive, it's unkind And passive aggressive communication means sarcasm
}

\oralblock{第一次关键诊断}{
  \begin{itemize}
    \item \textbf{内容}：已经抓住文章主旨，也主动说出了三类 communication styles；主要瓶颈不是文章完全没懂。
    \item \textbf{边界}：assertive 不能只压成 \textit{firm and direct}，必须保留 respectful / fair；passive-aggressive 也不等于单纯 sarcasm，它更广，核心是 indirect conflict + hidden negative message。
    \item \textbf{表达缺口}：\textit{how to distinct these} 应为 \textit{how to distinguish among these communication styles}。该主动缺口建立为 \texttt{EXP-012}。
    \item \textbf{文字层}：\textit{There are many nuances in English...} 比 \textit{There is many nuance...} 自然。
    \item 只有转写，没有原始音频，因此不据此评价发音、停顿、语速或真实流利度。
  \end{itemize}
}

\oralblock{AI 参考复述}{
  English communication can be more nuanced than the literal words suggest. This article teaches readers to distinguish among assertive, aggressive, and passive-aggressive styles and to read between the lines. Assertive communication is firm but fair: people express their needs clearly while still respecting others. Aggressive communication is also direct, but it can be rude and cause unnecessary upset. Passive-aggressive communication avoids direct conflict and often uses sarcasm or says the opposite of what the speaker really means. The quiz shows that a polite-sounding sentence can sometimes hide criticism or annoyance.
}

\oralblock{对照：第一次输出 vs 参考版}{
  \begin{itemize}
    \item \textbf{已经说出来的}：nuance、文章目标、三种沟通方式、aggressive 的负面性、passive-aggressive 与 sarcasm 的联系。
    \item \textbf{必须补回的语义边界}：assertive = respectful / fair；passive-aggressive = indirect / hidden negative message。
    \item \textbf{值得带走的语块}：\textit{distinguish among...}；\textit{firm but fair}；\textit{read between the lines}；\textit{avoid direct conflict}。
    \item \textbf{结构}：文章目标 → 三类定义 → quiz 用真实场景检验，不需要复述十道题的每个细节。
  \end{itemize}
}

\oralblock{即时回看后的第二次复述 · 原始转写}{
  Okay, I will repeat what I learned. This passage mainly introduces the nuance of English communication. The article teaches readers to distinguish among assertive, aggressive, and passive aggressive communication. A authoritative communication means firm and fair. This expression I made first and I can't get what it means. And now I know it means express their thoughts directly and respect others as well. As for aggressive communication, it often lead to absence, and passive aggressive often in polite expression and like sarcasm. Also, the article offers several conversations and teach me how to read the air.
}

\oralblock{第二次复述后的最小反馈}{
  \begin{itemize}
    \item \textbf{已经修复}：能正确调用 \textit{distinguish among...}，也主动补回 assertive 的 respectful 边界；由于刚看过参考版，这只支持 \texttt{EXP-012} 到 \texttt{L2 · 提示使用}。
    \item \textbf{概念词}：应为 \textit{assertive communication}，不是 \textit{authoritative communication}；原文固定概括是 \textit{firm but fair}。
    \item \textbf{aggressive}：转写中的 \textit{absence} 不符合原文；原文是 \textit{cause unnecessary upset}。若是 ASR 误转，不据此判断发音。
    \item \textbf{passive-aggressive}：已经从“= sarcasm”进步到“表面可能礼貌 + 实际带负面意味”；完整边界仍是 indirect + hidden negative message。
    \item \textbf{文章任务}：本文表达是 \textit{read between the lines}，不是 \textit{read the air}。
  \end{itemize}
}

\oralblock{本篇最终可达到版 · 不要求当天继续背诵}{
  This passage mainly introduces the nuances of English communication. The article teaches readers to distinguish among assertive, aggressive, and passive-aggressive communication. Assertive communication means being firm but fair: people express their thoughts directly while still respecting others. Aggressive communication is also direct, but it can cause unnecessary upset. Passive-aggressive communication often sounds polite on the surface, but it can express criticism or annoyance through sarcasm or hidden messages. The article also gives several conversations to teach readers how to read between the lines.
}

% ------------------------------------------------------------------------------
% 五、本篇收束：今天到此为止，不继续同篇重复
% ------------------------------------------------------------------------------
\readingsection{五、本篇收束与跨层引用}

\begin{itemize}
  \item \textbf{首读标记已闭环}：P1 两个 [!]、P3 [?]、P13 [?]、P16 [!]、P29 [?]，以及 P27 额外疑问均已处理。
  \item \textbf{Quiz 已校准}：隐藏攻击项为 1C、2A、3C、4B、5A、6B、7A、8B、9C、10A；重点不是背答案，而是识别“表面措辞”和“真实意图”的不一致。
  \item \textbf{进入表达层}：\texttt{EXP-012} · \textit{distinguish A from B / distinguish among...}，当前为 \texttt{L2 · 提示使用}。
  \item \textbf{仍留本篇}：\textit{firm but fair}、\textit{read between the lines}、\textit{put someone on their guard}、\textit{take something literally} 等先作为本篇高价值语块，不为凑数量全部建长期节点。
  \item \textbf{能力证据}：第二篇文章再次说明，当前理解内容的能力高于即时英语压缩能力；当需要同时回忆定义和组织英语时，容易把必要语义边界压掉。相关结论已在能力层记录。
  \item \textbf{今天停止条件已满足}：已经完成独立首读、标记闭环、quiz 校准、真实首答、参考版对照和一次即时回看。\textbf{今天不再要求整篇复述。}
\end{itemize}

\begin{coretakeaway}[下一次只做一次短冷检索]
  隔开一个训练窗口后，先不看本页，只回答：\textit{What is the difference between assertive, aggressive, and passive-aggressive communication?} 如果能自然带出 \textit{distinguish among...} 或 \textit{read between the lines}，再记录新的无提示证据；否则只局部回看，不重新走整篇流程。
\end{coretakeaway}

\end{document}
